%!TEX program = lualatex
\documentclass[a5paper,ngerman,headheight=80pt,10pt,DIV=calc]{scrartcl}
\usepackage[ngerman]{babel}
\usepackage{blindtext}
\usepackage[utf8]{luainputenc}
\usepackage{scrlayer-scrpage}
\usepackage[T1]{fontenc}
\usepackage{fontspec}
\usepackage{multicol}
\usepackage{enumitem}
\usepackage{epigraph}
\usepackage[german=quotes]{csquotes}
\usepackage[footskip=20pt,bottom=50pt]{geometry}
\setmainfont{Source Sans Pro}
\MakeOuterQuote{"}

\setlength{\parindent}{0pt}
\setlength{\parskip}{4pt}

\pagestyle{plain.scrheadings}
\thispagestyle{scrheadings}
\lohead{\vspace*{-1.1cm}\includegraphics[scale=0.3]{../images/stupa_logo_rot.pdf}}
\rohead{\\Präsidium des\\Studierendenparlaments}

\lofoot[]{V.i.S.d.P. Gunhild Berdal, Jim Martens, Ramon Weilinger\\Von-Melle-Park 5, 20146 Hamburg}
\cofoot[]{}
\cohead[\thepage]{}

\renewcommand{\thesubsection}{TOP \arabic{subsection}}
\renewcommand{\thesubsubsection}{\arabic{subsubsection}.}
\setcounter{subsection}{-1}

\begin{document}
    UHH · StuPa-Präsidium · Von-Melle-Park 5 · D-20146 Hamburg\\
    Tel.: 040/450 204 39 · stupa@uni-hamburg.de

    \subsection*{Einladung zur Sitzung des Studierendenparlaments am 30. Juni}
    \subsubsection*{Diskussion mit Wissenschaftssenatorin Katharina Fegebank}

    Bereits zu ihrem Antritt als Zweite Bürgermeisterin Hamburgs und Senatorin
    für Wissenschaft, Forschung und Gleichstellung hat Katharina Fegebank eine
    hohe Kommunikationsbereitschaft gegenüber den Hochschulen signalisiert. Doch
    die Gestaltung von Wissenschaftspolitik vollzieht sich heute in einem
    großen Spannungsfeld: Auf der einen Seite stehen der gesellschaftliche
    Auftrag der Wissenschaft, den globalen Herausforderungen gerecht zu werden,
    sowie zur Bildung mündiger Menschen beizutragen. Auf der anderen Seite stehen
    die chronische Unterfinanzierung der Hochschulen und der wachsende
    Konkurrenzdruck.

    Wir haben die Wissenschaftssenatorin eingeladen, um in dieser bisweilen
    konfliktreichen Lage mit den Studierenden zu diskutieren, was die jeweiligen
    Vorstellungen und Ansprüche für die Entwicklung der Wissenschaft(en) sind.
    Die Themen, die uns hauptsächlich beschäftigen werden, sind die
    Exzellenzstrategie, die Finanzierung der Hochschulen und die mögliche Schaffung
    einer "Zivilklausel" für die Forschung.

    Ihr seid herzlich eingeladen an der Diskussion teilzunehmen.
    \textbf{Die Sitzung ist öffentlich und alle haben Rederecht.}

    \begin{center}
        \textbf{
        Donnerstag, 30. Juni 2016,\\
        19-20 Uhr\\
        im Raum S28 der ehemaligen HWP\\
        (Von-Melle-Park 9, 20146 Hamburg)}
    \end{center}
    Auf der Tagesordnung der StuPa-Sitzung stehen danach u.a. die Verabschiedung
    des Haushalts der Verfassten Studierendenschaft, ein Antrag zur Verlegung
    der StuPa-Wahl ins Sommersemester sowie weitere aktuelle Themen.
\end{document}
